\documentclass[11pt]{article}
\usepackage[spanish,es-tabla,es-nodecimaldot]{babel}
\usepackage[utf8]{inputenc}
\usepackage[T1]{fontenc}
\usepackage[margin=2.2cm]{geometry}
\usepackage{amsmath,amssymb}
\usepackage{booktabs}
\usepackage{array}
\usepackage{xcolor}
\usepackage{mathtools}

\newcommand{\mm}[1]{\mathbf{#1}}

\title{\vspace{-1.5em} Ejercicio Regresión L\'ineal M\'ultiple}
\author{}
\date{}

\begin{document}
\maketitle
\vspace{-2.5em}

\section{Modelo y datos}

Modelo de regresión lineal múltiple con $n=5$ individuos:
\[
  y_i = \beta_0 + \beta_1 x_{i1} + \beta_2 x_{i2} + u_i, \qquad i=1,\dots,5,
\]
donde $x_{1}=$ experiencia (años) y $x_{2}=$ educación (años).

\begin{center}
\begin{tabular}{lcccc}
\toprule
Individuo & \emph{unos} & $x_1$ (experiencia) & $x_2$ (educación) & $y$ (salario) \\
\midrule
Juan  & 1 & 12 & 93  & 769 \\
Pedro & 1 & 18 & 119 & 808 \\
Sofía & 1 & 14 & 108 & 825 \\
Caro  & 1 & 12 & 96  & 650 \\
Alex  & 1 & 11 & 74  & 562 \\
\bottomrule
\end{tabular}
\end{center}

\bigskip
\noindent Matrices del modelo:
\[
  X=\begin{bmatrix}
    1 & 12 & 93\\
    1 & 18 & 119\\
    1 & 14 & 108\\
    1 & 12 & 96\\
    1 & 11 & 74
  \end{bmatrix},
  \qquad
  X^\top=\begin{bmatrix}
    1 & 1 & 1 & 1 & 1\\
    12 & 18 & 14 & 12 & 11\\
    93 & 119 & 108 & 96 & 74
  \end{bmatrix}.
\]

Multiplicando $X^\top X$ (cada entrada es la suma de productos correspondiente: $n$, $\sum x_1$, $\sum x_2$, $\sum x_1^2$, $\sum x_1 x_2$, $\sum x_2^2$):
\[
  A \;=\; X^\top X \;=\;
  \begin{bmatrix}
    5 & 67 & 490\\
    67 & 929 & 6.736\\
    490 & 6.736 & 49.166
  \end{bmatrix}.
\]

El objetivo de este documento es obtener $A^{-1}=(X^\top X)^{-1}$ \textbf{paso a paso, usando el método de cofactores}:
\[
  A^{-1} = \frac{1}{\det(A)}\,\operatorname{adj}(A),
  \qquad
  \operatorname{adj}(A) = C^\top,
\]
donde $C$ es la matriz de cofactores de $A$.

\section{Paso 1: Matriz de menores}

El menor $M_{ij}$ es el determinante de la submatriz $2\times 2$ que resulta de eliminar la fila $i$ y la columna $j$ de $A$.

\[
M_{11}=\begin{vmatrix}929 & 6.736\\ 6.736 & 49.166\end{vmatrix}
= 929\times 49.166 - 6.736\times 6.736
= 45.675.214-45.373.696=301.518
\]
\[
M_{12}=\begin{vmatrix}67 & 6.736\\ 490 & 49.166\end{vmatrix}
= 67\times 49.166 - 6.736\times 490
= 3.294.122-3.300.640=-6.518
\]
\[
M_{13}=\begin{vmatrix}67 & 929\\ 490 & 6.736\end{vmatrix}
= 67\times 6.736 - 929\times 490
= 451.312-455.210=-3.898
\]
\[
M_{21}=\begin{vmatrix}67 & 490\\ 6.736 & 49.166\end{vmatrix}
= 67\times 49.166 - 490\times 6.736
= 3.294.122-3.300.640=-6.518
\]
\[
M_{22}=\begin{vmatrix}5 & 490\\ 490 & 49.166\end{vmatrix}
= 5\times 49.166 - 490\times 490
= 245.830-240.100=5.730
\]
\[
M_{23}=\begin{vmatrix}5 & 67\\ 490 & 6.736\end{vmatrix}
= 5\times 6.736 - 67\times 490
= 33.680-32.830=850
\]
\[
M_{31}=\begin{vmatrix}67 & 490\\ 929 & 6.736\end{vmatrix}
= 67\times 6.736 - 490\times 929
= 451.312-455.210=-3.898
\]
\[
M_{32}=\begin{vmatrix}5 & 490\\ 67 & 6.736\end{vmatrix}
= 5\times 6.736 - 490\times 67
= 33.680-32.830=850
\]
\[
M_{33}=\begin{vmatrix}5 & 67\\ 67 & 929\end{vmatrix}
= 5\times 929 - 67\times 67
= 4.645-4.489=156
\]

Por lo tanto, la matriz de menores es:
\[
  M=\begin{bmatrix}
    301.518 & -6.518 & -3.898\\
    -6.518 & 5.730 & 850\\
    -3.898 & 850 & 156
  \end{bmatrix}.
\]

\emph{(Nota: como $A=X^\top X$ es simétrica, se cumple $M_{12}=M_{21}$ y $M_{13}=M_{31}$, $M_{23}=M_{32}$, lo cual sirve como chequeo intermedio.)}

\section{Paso 2: Matriz de cofactores}

Cada cofactor se obtiene aplicando el patrón de signos (tablero de ajedrez) a los menores:
\[
  \operatorname{sgn}(i,j)=(-1)^{i+j}
  \;\Longrightarrow\;
  \begin{bmatrix}
    + & - & +\\
    - & + & -\\
    + & - & +
  \end{bmatrix}.
\]

\[
  C_{11}=+M_{11}=301.518, \quad
  C_{12}=-M_{12}=6.518, \quad
  C_{13}=+M_{13}=-3.898,
\]
\[
  C_{21}=-M_{21}=6.518, \quad
  C_{22}=+M_{22}=5.730, \quad
  C_{23}=-M_{23}=-850,
\]
\[
  C_{31}=+M_{31}=-3.898, \quad
  C_{32}=-M_{32}=-850, \quad
  C_{33}=+M_{33}=156.
\]

\[
  C=\begin{bmatrix}
    301.518 & 6.518 & -3.898\\
    6.518 & 5.730 & -850\\
    -3.898 & -850 & 156
  \end{bmatrix}.
\]

\section{Paso 3: Matriz adjunta (adjunta = transpuesta de cofactores)}

\[
  \operatorname{adj}(A) = C^\top
  =\begin{bmatrix}
    301.518 & 6.518 & -3.898\\
    6.518 & 5.730 & -850\\
    -3.898 & -850 & 156
  \end{bmatrix}.
\]

Como $A$ es simétrica, $C$ también resulta simétrica, así que $\operatorname{adj}(A)=C$ en este caso particular (esto \textbf{no} ocurre en general; solo porque $A=X^\top X$ es simétrica).

\section{Paso 4: Determinante de $A$}

El determinante se calcula expandiendo por cofactores a lo largo de la primera fila:
\[
  \det(A) = a_{11}C_{11}+a_{12}C_{12}+a_{13}C_{13}
\]
\[
  \det(A) = 5\times 301.518 \;+\; 67\times 6.518 \;+\; 490\times(-3.898)
\]
\[
  \det(A) = 1.507.590 \;+\; 436.706 \;-\; 1.910.020
\]
\[
  \boxed{\det(A) = 34.276}
\]

\section{Paso 5: Matriz inversa}

\[
  A^{-1} = \frac{1}{\det(A)}\operatorname{adj}(A)
  = \frac{1}{34.276}
  \begin{bmatrix}
    301.518 & 6.518 & -3.898\\
    6.518 & 5.730 & -850\\
    -3.898 & -850 & 156
  \end{bmatrix}.
\]

Dividiendo cada entrada por $34.276$ se obtiene, en decimales:

\[
  A^{-1} = (X^\top X)^{-1} \approx
  \begin{bmatrix}
    8,796767 & 0,190162 & -0,113724\\
    0,190162 & 0,167172 & -0,024799\\
    -0,113724 & -0,024799 & 0,004551
  \end{bmatrix}.
\]


\section{ Cálculo de $\hat\beta=(X^\top X)^{-1}X^\top y$}

Como la base de datos ya incluye la variable $y$ (salario), es posible completar el ejercicio calculando el vector $X^\top y$ y luego el estimador OLS. Con $y=(769,\,808,\,825,\,650,\,562)^\top$:

\[
  X^\top y =
  \begin{bmatrix}
    \sum y_i \\ \sum x_{1i} y_i \\ \sum x_{2i} y_i
  \end{bmatrix}
  =
  \begin{bmatrix}
    3.614\\
    49.304\\
    360.757
  \end{bmatrix}.
\]

Luego:
\[
  \hat\beta = A^{-1}X^\top y \approx
  \begin{bmatrix}
    140,586\\
    -16,790\\
    8,237
  \end{bmatrix}.
\]

Es decir, $\hat\beta_0\approx 140,586$, $\hat\beta_1\approx -16,790$ y $\hat\beta_2\approx 8,237$.

\end{document}